\documentclass[11pt,a4paper]{article}
\usepackage[left=1.35cm,right=1.35cm,top=1.15cm,bottom=1.15cm]{geometry}
\usepackage{fontspec}
\usepackage{xcolor}
\usepackage{enumitem}
\usepackage{ragged2e}
\usepackage{titlesec}
\usepackage{hyperref}
\setmainfont{texgyreheros-regular.otf}[BoldFont=texgyreheros-bold.otf,ItalicFont=texgyreheros-italic.otf,BoldItalicFont=texgyreheros-bolditalic.otf]
\definecolor{navy}{HTML}{17365D}
\definecolor{linkblue}{HTML}{1F5A94}
\hypersetup{colorlinks=true,urlcolor=linkblue,linkcolor=linkblue,pdfauthor={Qixuan Yang},pdftitle={Qixuan Yang - Software Engineering Group Projects CV}}
\pagestyle{empty}
\setlength{\parindent}{0pt}
\setlength{\parskip}{1pt}
\linespread{1.02}
\setlength{\emergencystretch}{1.5em}
\setlist[itemize]{leftmargin=1.35em,itemsep=2pt,topsep=2pt,parsep=0pt,before=\RaggedRight}
\titleformat{\section}{\large\bfseries\color{navy}}{}{0pt}{}[\vspace{-3pt}\color{navy}\titlerule]
\titlespacing*{\section}{0pt}{8pt}{4pt}
\newcommand{\project}[1]{\par\vspace{2pt}\textbf{#1}\par}
% Standalone XeLaTeX/Tectonic source for Pacman, AI Game Player and Crisis App.
% Past experience is supported by the internship summary and existing CVs.
% Profile interests are intended contributions, not completed game projects.
\begin{document}
\begin{center}
{\LARGE\bfseries\color{navy} Qixuan Yang}\\[3pt]
{\small\href{mailto:scyqy5@nottingham.ac.uk}{scyqy5@nottingham.ac.uk}\enspace|\enspace\href{mailto:mark44yang@outlook.com}{mark44yang@outlook.com}\\[2pt]
+86 13306576505\enspace|\enspace\href{https://github.com/MarkYang44}{github.com/MarkYang44}}
\end{center}

\section*{Education}
\textbf{University of Nottingham, UK} \hfill Sep 2026 - Jun 2028 (expected)\\
BSc Computer Science with Artificial Intelligence (2+2), Jubilee Campus.\\[2pt]
\textbf{University of Nottingham Ningbo China} \hfill 2024 - 2026\\
First two years of the 2+2 program; GPA: 4.0/4.0 (Top 10\%).\\
Dean's Scholarship: 2025 - 2026.\quad IELTS: 7.5.

\section*{Profile}
Computer Science and AI student with hands-on Java, Python and Go development, AI prototyping and system integration experience. Interested in applying modular design, reliable state handling and model evaluation to games, AI opponents and user-facing reporting applications.

\section*{Relevant Experience}
% Official role category: AI类. Division names are descriptive translations.
% Level 1: 大数据应用事业部; Level 2: 数据要素事业部.
\textbf{China Telecom Fufu Information Technology Co., Ltd.}\hfill Jul 2026 - Aug 2026\\
AI Intern\enspace|\enspace Big Data Applications Division / Data Elements Division

\project{Factory safety AI assistant and AI workflows}
\begin{itemize}
\item Independently built a Python/Streamlit safety-analysis demo combining rules, RAG and DeepSeek to identify hazards, cite regulations and prioritize corrective actions; used ChromaDB and SQLite for retrieval and traceable task history.
\item Separated analysis services from the UI and handled retrieval and task failures. Built a LangGraph plan-and-execute agent prototype with tool routing and state management.
\end{itemize}

\project{Healthcare portal backend}
\begin{itemize}
\item Developed Java/Spring Boot REST APIs for homepage configuration, business metrics and news; implemented MyBatis/MySQL data aggregation, Nacos configuration and error handling.
\item Clarified database and gateway dependencies with frontend and deployment colleagues and prepared integration and release documentation, supporting coordinated delivery within an existing microservice system.
\end{itemize}

\project{Hyperledger Fabric platform and Go chaincode}
\begin{itemize}
\item Fixed record-ordering and input-validation defects in Go chaincode; added regression coverage while preserving existing API and ledger compatibility.
\item Rebuilt the local platform with WSL/Linux and Docker Compose and automated startup and end-to-end checks, strengthening reproducible integration and verification.
\end{itemize}

\section*{Academic Project}
\textbf{Image Classification System}\enspace|\enspace 2025
\begin{itemize}
\item Built a modular PyTorch/ResNet18 transfer-learning pipeline for preprocessing, training and evaluation; applied data augmentation and compared optimizers through repeatable experiments.
\item Containerized a Flask REST inference service with Docker, separating model execution from the application interface for reproducible deployment.
\end{itemize}

\section*{Technical Skills}
\textbf{Programming:} Java, Python, Go, C/C++, SQL, Shell, MATLAB.\\
\textbf{AI \& data:} PyTorch, model evaluation, LangGraph/LangChain, RAG, ChromaDB, Pandas, SQLite.\\
\textbf{Backend \& interfaces:} Spring Boot, REST APIs, MyBatis-Plus, MySQL, Nacos, Flask, Streamlit.\\
\textbf{Engineering:} Git, Docker Compose, WSL/Linux, JUnit, Go testing, regression testing, integration documentation, Hyperledger Fabric.

\end{document}
